\chapter{Modelo Lógico e Dicionário de Dados (DER Lógico)}
\label{cap_modelo_logico}

O Modelo Lógico de Dados traduz as abstrações do modelo conceitual para as estruturas formais do modelo relacional \cite{codd1970, martin1989}, definindo a tipagem dos dados no Sistema Gerenciador de Banco de Dados (PostgreSQL 16 LTS), a integridade referencial por meio de chaves primárias e estrangeiras e as restrições canônicas de domínio.

% ===========================================================================
\section{Diagrama Entidade-Relacionamento Lógico (DER)}
\label{sec_diagrama_der_logico}
% ===========================================================================

A Figura \ref{fig_der_logico} apresenta o Diagrama Entidade-Relacionamento Lógico (DER) completo do SIGEA-GTT, modelado com as oito tabelas físicas relacionais e seus relacionamentos estruturais na notação Pé de Galinha (\textit{Crow's Foot}).

\begin{figure}[!ht]
\centering
\caption{Diagrama Entidade-Relacionamento Lógico (Notação Crow's Foot / Relacional)}
\label{fig_der_logico}
\resizebox{0.95\textwidth}{!}{%
\begin{tikzpicture}[
    node distance=1.5cm and 2cm,
    table/.style={
        rectangle split,
        rectangle split parts=2,
        rectangle split part fill={blue!20, white},
        draw=blue!80!black,
        thick,
        rounded corners=2pt,
        font=\scriptsize,
        minimum width=4.6cm,
        text width=4.4cm,
        align=left
    },
    table_aux/.style={
        rectangle split,
        rectangle split parts=2,
        rectangle split part fill={teal!20, white},
        draw=teal!80!black,
        thick,
        rounded corners=2pt,
        font=\scriptsize,
        minimum width=4.6cm,
        text width=4.4cm,
        align=left
    },
    fk_line/.style={
        draw=black!70,
        thick,
        -Latex
    }
]

    % ====================================================
    % TABELA USERS
    % ====================================================
    \node[table] (users) at (0, 7) {
        \textbf{\normalsize users}
        \nodepart{two}
        \textbf{PK} id : UUID\\
        \textbf{UK} email : VARCHAR(120)\\
        full\_name : VARCHAR(150)\\
        password\_hash : VARCHAR(255)\\
        role : user\_role\\
        registration\_number : VARCHAR(30)\\
        is\_active : BOOLEAN\\
        must\_change\_password : BOOLEAN\\
        created\_at : TIMESTAMPTZ\\
        updated\_at : TIMESTAMPTZ
    };

    % ====================================================
    % TABELA ACADEMIC_CLASSES
    % ====================================================
    \node[table] (classes) at (7.5, 7) {
        \textbf{\normalsize academic\_classes}
        \nodepart{two}
        \textbf{PK} id : UUID\\
        subject\_name : VARCHAR(120)\\
        class\_code : VARCHAR(20)\\
        academic\_period : VARCHAR(10)\\
        \textbf{FK} professor\_id : UUID\\
        is\_closed : BOOLEAN\\
        created\_at : TIMESTAMPTZ
    };

    % ====================================================
    % TABELA CLASS_STUDENTS
    % ====================================================
    \node[table_aux] (class_students) at (3.8, 2.5) {
        \textbf{\normalsize class\_students}
        \nodepart{two}
        \textbf{PK, FK1} class\_id : UUID\\
        \textbf{PK, FK2} student\_id : UUID\\
        enrolled\_at : TIMESTAMPTZ
    };

    % ====================================================
    % TABELA ACTIVITIES
    % ====================================================
    \node[table] (activities) at (12.5, 7) {
        \textbf{\normalsize activities}
        \nodepart{two}
        \textbf{PK} id : UUID\\
        \textbf{FK} class\_id : UUID\\
        title : VARCHAR(150)\\
        description : TEXT\\
        clinical\_case\_data : JSONB\\
        deadline : TIMESTAMPTZ\\
        created\_at : TIMESTAMPTZ
    };

    % ====================================================
    % TABELA ACTIVITY_SUBMISSIONS
    % ====================================================
    \node[table] (submissions) at (12.5, 1.8) {
        \textbf{\normalsize activity\_submissions}
        \nodepart{two}
        \textbf{PK} id : UUID\\
        \textbf{FK1} activity\_id : UUID\\
        \textbf{FK2} student\_id : UUID\\
        identified\_triggers : JSONB\\
        quality\_tools\_data : JSONB\\
        submission\_date : TIMESTAMPTZ\\
        grade : NUMERIC(4,2)\\
        professor\_feedback : TEXT\\
        graded\_at : TIMESTAMPTZ
    };

    % ====================================================
    % TABELA GTT_MODULES
    % ====================================================
    \node[table] (gtt_modules) at (0, -2.5) {
        \textbf{\normalsize gtt\_modules}
        \nodepart{two}
        \textbf{PK} id : UUID\\
        \textbf{UK} code : VARCHAR(10)\\
        name : VARCHAR(100)\\
        description : TEXT\\
        is\_active : BOOLEAN\\
        created\_at : TIMESTAMPTZ
    };

    % ====================================================
    % TABELA GTT_TRIGGERS
    % ====================================================
    \node[table] (gtt_triggers) at (6.5, -2.5) {
        \textbf{\normalsize gtt\_triggers}
        \nodepart{two}
        \textbf{PK} id : UUID\\
        \textbf{FK} module\_id : UUID\\
        \textbf{UK} code : VARCHAR(10)\\
        name : VARCHAR(150)\\
        description : TEXT\\
        is\_active : BOOLEAN\\
        created\_at : TIMESTAMPTZ
    };

    % ====================================================
    % TABELA HARM_SEVERITIES
    % ====================================================
    \node[table] (harm) at (12.5, -2.5) {
        \textbf{\normalsize harm\_severities}
        \nodepart{two}
        \textbf{PK} id : UUID\\
        \textbf{UK} category\_letter : VARCHAR(1)\\
        name : VARCHAR(100)\\
        description : TEXT\\
        is\_harm : BOOLEAN\\
        is\_active : BOOLEAN
    };

    % ====================================================
    % CONEXÕES DE CHAVES ESTRANGEIRAS
    % ====================================================
    % users -> academic_classes (professor_id)
    \draw[fk_line] (users.east) -- node[above, font=\tiny] {1 : N} (classes.west);

    % academic_classes -> class_students (class_id)
    \draw[fk_line] (classes.south) -- ++(0,-1) -| (class_students.north);

    % users -> class_students (student_id)
    \draw[fk_line] (users.south) -- ++(0,-1) -| (class_students.north);

    % academic_classes -> activities (class_id)
    \draw[fk_line] (classes.east) -- node[above, font=\tiny] {1 : N} (activities.west);

    % activities -> activity_submissions (activity_id)
    \draw[fk_line] (activities.south) -- node[right, font=\tiny] {1 : N} (submissions.north);

    % users -> activity_submissions (student_id)
    \draw[fk_line] (users.south) -- ++(0,-2.5) -| ++(9,0) |- (submissions.west);

    % gtt_modules -> gtt_triggers (module_id)
    \draw[fk_line] (gtt_modules.east) -- node[above, font=\tiny] {1 : N} (gtt_triggers.west);

\end{tikzpicture}
}
\fonte{Elaborado pelo autor (2026).}
\end{figure}

% ===========================================================================
\section{Dicionário de Dados das Tabelas Físicas}
\label{sec_dicionario_dados}
% ===========================================================================

A seguir são apresentadas as especificações detalhadas das oito tabelas implementadas no banco de dados do SIGEA-GTT.

\subsection{Tabela: \texttt{users}}
Centraliza os registros de credenciais e permissões de todos os atores do sistema.

\begin{table}[!ht]
\centering
\caption{Dicionário de Dados da Tabela \texttt{users}}
\label{tab_dic_users}
\footnotesize
\begin{tabularx}{\textwidth}{p{3.2cm} p{2.5cm} c c p{1.8cm} X}
\toprule
\textbf{Atributo} & \textbf{Tipo} & \textbf{Nulo?} & \textbf{Chave} & \textbf{Padrão} & \textbf{Descrição e Regra de Negócio} \\
\midrule
\texttt{id} & UUID & Não & PK & \texttt{uuid\_v4()} & Identificador universal único do usuário. \\
\texttt{full\_name} & VARCHAR(150) & Não & -- & -- & Nome civil completo do usuário. \\
\texttt{email} & VARCHAR(120) & Não & UK & -- & E-mail institucional exclusivo (\texttt{@academico.ufs.br}). \\
\texttt{password\_hash} & VARCHAR(255) & Não & -- & -- & Resumo criptográfico da senha (BCrypt). \\
\texttt{role} & \texttt{user\_role} & Não & -- & -- & Perfil de acesso: \texttt{'ADMIN'}, \texttt{'PROFESSOR'} ou \texttt{'STUDENT'}. \\
\texttt{registration\_\-number} & VARCHAR(30) & Sim & -- & NULL & Matrícula acadêmica. Obrigatória para discente e nula para os demais. \\
\texttt{is\_active} & BOOLEAN & Não & -- & \texttt{TRUE} & Situação cadastral do usuário (ativo/inativo). \\
\texttt{must\_\-change\_\-password} & BOOLEAN & Não & -- & \texttt{TRUE} & Flag de primeiro acesso para forçamento de troca de credencial. \\
\texttt{created\_at} & TIMESTAMPTZ & Não & -- & \texttt{NOW()} & Carimbo de data/hora de criação do registro. \\
\texttt{updated\_at} & TIMESTAMPTZ & Não & -- & \texttt{NOW()} & Carimbo de data/hora da última alteração de perfil/senha. \\
\bottomrule
\end{tabularx}
\fonte{Elaborado pelo autor (2026).}
\end{table}

\subsection{Tabela: \texttt{gtt\_modules}}
Armazena os seis módulos canônicos de rastreamento do \textit{Institute for Healthcare Improvement}.

\begin{table}[!ht]
\centering
\caption{Dicionário de Dados da Tabela \texttt{gtt\_modules}}
\label{tab_dic_gtt_modules}
\footnotesize
\begin{tabularx}{\textwidth}{p{3.2cm} p{2.5cm} c c p{1.8cm} X}
\toprule
\textbf{Atributo} & \textbf{Tipo} & \textbf{Nulo?} & \textbf{Chave} & \textbf{Padrão} & \textbf{Descrição e Regra de Negócio} \\
\midrule
\texttt{id} & UUID & Não & PK & \texttt{uuid\_v4()} & Identificador universal único do módulo. \\
\texttt{code} & VARCHAR(10) & Não & UK & -- & Código alfanumérico mnemônico internacional (\texttt{C}, \texttt{S}, \texttt{M}, \texttt{I}, \texttt{P}, \texttt{E}). \\
\texttt{name} & VARCHAR(100) & Não & -- & -- & Denominação oficial do módulo clínico do IHI. \\
\texttt{description} & TEXT & Sim & -- & NULL & Escopo assistencial e instruções de busca do módulo. \\
\texttt{is\_active} & BOOLEAN & Não & -- & \texttt{TRUE} & Status de ativação curricular do módulo. \\
\texttt{created\_at} & TIMESTAMPTZ & Não & -- & \texttt{NOW()} & Data e hora de inclusão do módulo no sistema. \\
\bottomrule
\end{tabularx}
\fonte{Elaborado pelo autor (2026).}
\end{table}

\subsection{Tabela: \texttt{gtt\_triggers}}
Contém os 53 gatilhos clínicos padronizados vinculados aos módulos de auditoria.

\begin{table}[!ht]
\centering
\caption{Dicionário de Dados da Tabela \texttt{gtt\_triggers}}
\label{tab_dic_gtt_triggers}
\footnotesize
\begin{tabularx}{\textwidth}{p{3.2cm} p{2.5cm} c c p{1.8cm} X}
\toprule
\textbf{Atributo} & \textbf{Tipo} & \textbf{Nulo?} & \textbf{Chave} & \textbf{Padrão} & \textbf{Descrição e Regra de Negócio} \\
\midrule
\texttt{id} & UUID & Não & PK & \texttt{uuid\_v4()} & Identificador universal único do gatilho. \\
\texttt{module\_id} & UUID & Não & FK & -- & Referência ao módulo (\texttt{gtt\_modules.id}) com \texttt{RESTRICT}. \\
\texttt{code} & VARCHAR(10) & Não & UK & -- & Código único do gatilho (ex.: \texttt{'C01'}, \texttt{'M03'}, \texttt{'S08'}). \\
\texttt{name} & VARCHAR(150) & Não & -- & -- & Nome clínico consagrado do gatilho rastreador. \\
\texttt{description} & TEXT & Não & -- & -- & Diretrizes operacionais e critérios de detecção em prontuário. \\
\texttt{is\_active} & BOOLEAN & Não & -- & \texttt{TRUE} & Indicador de uso do gatilho em atividades ativas. \\
\texttt{created\_at} & TIMESTAMPTZ & Não & -- & \texttt{NOW()} & Data e hora de cadastro do gatilho. \\
\bottomrule
\end{tabularx}
\fonte{Elaborado pelo autor (2026).}
\end{table}

\subsection{Tabela: \texttt{harm\_severities}}
Registra as nove categorias de gravidade de dano estabelecidas pelo conselho NCC MERP.

\begin{table}[!ht]
\centering
\caption{Dicionário de Dados da Tabela \texttt{harm\_severities}}
\label{tab_dic_harm_severities}
\footnotesize
\begin{tabularx}{\textwidth}{p{3.2cm} p{2.5cm} c c p{1.8cm} X}
\toprule
\textbf{Atributo} & \textbf{Tipo} & \textbf{Nulo?} & \textbf{Chave} & \textbf{Padrão} & \textbf{Descrição e Regra de Negócio} \\
\midrule
\texttt{id} & UUID & Não & PK & \texttt{uuid\_v4()} & Identificador universal único da gravidade. \\
\texttt{category\_\-letter} & VARCHAR(1) & Não & UK & -- & Letra da categoria na escala NCC MERP (\texttt{'A'} até \texttt{'I'}). \\
\texttt{name} & VARCHAR(100) & Não & -- & -- & Denominação conceitual do grau de severidade. \\
\texttt{description} & TEXT & Não & -- & -- & Descrição das consequências clínicas e intervenções exigidas. \\
\texttt{is\_harm} & BOOLEAN & Não & -- & -- & \texttt{FALSE} para categorias \texttt{A}--\texttt{D}; \texttt{TRUE} para \texttt{E}--\texttt{I} (dano real). \\
\texttt{is\_active} & BOOLEAN & Não & -- & \texttt{TRUE} & Status de vigência da gravidade no sistema. \\
\bottomrule
\end{tabularx}
\fonte{Elaborado pelo autor (2026).}
\end{table}

\subsection{Tabela: \texttt{academic\_classes}}
Representa as turmas acadêmicas semestrais criadas sob responsabilidade docente.

\begin{table}[!ht]
\centering
\caption{Dicionário de Dados da Tabela \texttt{academic\_classes}}
\label{tab_dic_academic_classes}
\footnotesize
\begin{tabularx}{\textwidth}{p{3.2cm} p{2.5cm} c c p{1.8cm} X}
\toprule
\textbf{Atributo} & \textbf{Tipo} & \textbf{Nulo?} & \textbf{Chave} & \textbf{Padrão} & \textbf{Descrição e Regra de Negócio} \\
\midrule
\texttt{id} & UUID & Não & PK & \texttt{uuid\_v4()} & Identificador universal único da turma. \\
\texttt{subject\_name} & VARCHAR(120) & Não & -- & -- & Nome oficial do componente curricular universitário. \\
\texttt{class\_code} & VARCHAR(20) & Não & -- & -- & Código identificador da turma (ex.: \texttt{'T01'}, \texttt{'T02'}). \\
\texttt{academic\_period} & VARCHAR(10) & Não & -- & -- & Período letivo no formato semestral (ex.: \texttt{'2025.2'}, \texttt{'2026.1'}). \\
\texttt{professor\_id} & UUID & Não & FK & -- & Referência ao docente titular (\texttt{users.id}) com \texttt{RESTRICT}. \\
\texttt{is\_closed} & BOOLEAN & Não & -- & \texttt{FALSE} & Situação da turma (\texttt{TRUE} = encerrada e imutável; \texttt{FALSE} = ativa). \\
\texttt{created\_at} & TIMESTAMPTZ & Não & -- & \texttt{NOW()} & Data e hora de homologação da turma. \\
\bottomrule
\end{tabularx}
\fonte{Elaborado pelo autor (2026).}
\end{table}

\subsection{Tabela: \texttt{class\_students}}
Modela a relação muitos-para-muitos entre turmas e discentes inscritos.

\begin{table}[!ht]
\centering
\caption{Dicionário de Dados da Tabela \texttt{class\_students}}
\label{tab_dic_class_students}
\footnotesize
\begin{tabularx}{\textwidth}{p{3.2cm} p{2.5cm} c c p{1.8cm} X}
\toprule
\textbf{Atributo} & \textbf{Tipo} & \textbf{Nulo?} & \textbf{Chave} & \textbf{Padrão} & \textbf{Descrição e Regra de Negócio} \\
\midrule
\texttt{class\_id} & UUID & Não & PK, FK & -- & Referência à turma (\texttt{academic\_classes.id}) com \texttt{CASCADE}. \\
\texttt{student\_id} & UUID & Não & PK, FK & -- & Referência ao discente (\texttt{users.id}) com \texttt{RESTRICT}. \\
\texttt{enrolled\_at} & TIMESTAMPTZ & Não & -- & \texttt{NOW()} & Carimbo de data/hora da efetivação da inscrição acadêmica. \\
\bottomrule
\end{tabularx}
\fonte{Elaborado pelo autor (2026).}
\end{table}

\subsection{Tabela: \texttt{activities}}
Armazena as atividades de auditoria propostas pelo docente contendo o prontuário simulado.

\begin{table}[!ht]
\centering
\caption{Dicionário de Dados da Tabela \texttt{activities}}
\label{tab_dic_activities}
\footnotesize
\begin{tabularx}{\textwidth}{p{3.2cm} p{2.5cm} c c p{1.8cm} X}
\toprule
\textbf{Atributo} & \textbf{Tipo} & \textbf{Nulo?} & \textbf{Chave} & \textbf{Padrão} & \textbf{Descrição e Regra de Negócio} \\
\midrule
\texttt{id} & UUID & Não & PK & \texttt{uuid\_v4()} & Identificador universal único da atividade. \\
\texttt{class\_id} & UUID & Não & FK & -- & Referência à turma acadêmica (\texttt{academic\_classes.id}) com \texttt{CASCADE}. \\
\texttt{title} & VARCHAR(150) & Não & -- & -- & Título descritivo da prática pedagógica de auditoria. \\
\texttt{description} & TEXT & Não & -- & -- & Instruções e metas pedagógicas fornecidas pelo docente. \\
\texttt{clinical\_\-case\_\-data} & JSONB & Não & -- & -- & Prontuário simulado completo em formato documental semiestruturado. \\
\texttt{deadline} & TIMESTAMPTZ & Não & -- & -- & Data e hora limite para submissão das resoluções pelos alunos. \\
\texttt{created\_at} & TIMESTAMPTZ & Não & -- & \texttt{NOW()} & Data e hora de publicação da atividade. \\
\bottomrule
\end{tabularx}
\fonte{Elaborado pelo autor (2026).}
\end{table}

\subsection{Tabela: \texttt{activity\_submissions}}
Registra as auditorias concluídas pelos discentes com seus respectivos artefatos da qualidade e notas.

\begin{table}[!ht]
\centering
\caption{Dicionário de Dados da Tabela \texttt{activity\_submissions}}
\label{tab_dic_activity_submissions}
\footnotesize
\begin{tabularx}{\textwidth}{p{3.2cm} p{2.5cm} c c p{1.8cm} X}
\toprule
\textbf{Atributo} & \textbf{Tipo} & \textbf{Nulo?} & \textbf{Chave} & \textbf{Padrão} & \textbf{Descrição e Regra de Negócio} \\
\midrule
\texttt{id} & UUID & Não & PK & \texttt{uuid\_v4()} & Identificador universal único da submissão. \\
\texttt{activity\_id} & UUID & Não & FK & -- & Referência à atividade (\texttt{activities.id}) com \texttt{CASCADE}. \\
\texttt{student\_id} & UUID & Não & FK & -- & Referência ao discente (\texttt{users.id}) com \texttt{RESTRICT}. \\
\texttt{identified\_\-triggers} & JSONB & Não & -- & -- & Array documental de gatilhos detectados, gravidades e justificativas. \\
\texttt{quality\_\-tools\_\-data} & JSONB & Não & -- & -- & Dados estruturados das ferramentas: Ishikawa, 5W2H, GUT e PDCA. \\
\texttt{submission\_date} & TIMESTAMPTZ & Não & -- & \texttt{NOW()} & Carimbo temporal de submissão do trabalho pelo estudante. \\
\texttt{grade} & NUMERIC(4,2) & Sim & -- & NULL & Nota atribuída pelo docente entre \texttt{0.00} e \texttt{10.00}. \\
\texttt{professor\_\-feedback} & TEXT & Sim & -- & NULL & Parecer pedagógico qualitativo do professor titular. \\
\texttt{graded\_at} & TIMESTAMPTZ & Sim & -- & NULL & Carimbo de data/hora em que a avaliação foi concluída pelo docente. \\
\bottomrule
\end{tabularx}
\fonte{Elaborado pelo autor (2026).}
\end{table}
